% --- ОБЩИЕ НАСТРОЙКИ ТЕКСТА ---
\onehalfspacing
\setlength{\parindent}{1.25cm} % Абзацный отступ

% Улучшенное выравнивание по ширине
\sloppy 
\hyphenpenalty=1000 
\tolerance=2000

% --- ФИКС РАЗМЕРОВ ЗАГОЛОВКОВ (Убираем ОГРОМНЫЙ текст) ---
% disposition — это база для всех заголовков
\setkomafont{disposition}{\rmfamily\bfseries\normalsize}
\setkomafont{chapter}{\bfseries\normalsize}
\setkomafont{section}{\bfseries\normalsize}
\setkomafont{subsection}{\bfseries\normalsize}

% Настройка оглавления (чтобы само слово СОДЕРЖАНИЕ не было гигантским)
\setkomafont{chapterentry}{\normalfont\normalsize}
\setkomafont{chapterentrypagenumber}{\normalfont\normalsize}

% --- НАСТРОЙКА ОГЛАВЛЕНИЯ ---
\renewcaptionname{russian}{\contentsname}{\bfseries СОДЕРЖАНИЕ}
\BeforeTOCHead[toc]{
    \renewcommand*{\raggedsection}{\centering}
    \setkomafont{chapter}{\bfseries\normalsize} % Принудительно для заголовка TOC
}

% --- НАСТРОЙКА ОТСТУПОВ (1.25 СМ) ---
\KOMAoptions{numbers=noendperiod}

% Настраиваем отступы для всех уровней, включая главы
\RedeclareSectionCommand[
    indent=1.25cm, 
    beforeskip=-1sp, % Минимальный отступ сверху
    afterskip=6pt
]{chapter}

\RedeclareSectionCommand[
    indent=1.25cm, 
    beforeskip=12pt, 
    afterskip=6pt
]{section}

\RedeclareSectionCommand[
    indent=1.25cm, 
    beforeskip=12pt, 
    afterskip=6pt
]{subsection}

% --- СПИСКИ (enumitem) ---
\setlist{
	leftmargin=0pt,
	nosep,
	wide=\parindent,
	labelsep=0.5em,
	listparindent=\parindent
}
\setlist[itemize,1]{label=--} 
\setlist[enumerate,1]{label=\arabic*.}

% --- ПОДПИСИ И ТАБЛИЦЫ ---
\captionsetup{justification=centering, singlelinecheck=false}
\DeclareCaptionLabelFormat{simple}{#1~#2}
\captionsetup{labelformat=simple}
\DeclareCaptionLabelSeparator{emdash}{ --- }
\captionsetup[table]{
	labelsep=emdash,
	justification=raggedright,
	singlelinecheck=false,
	position=top,
	skip=6pt
}

% --- ЛИСТИНГИ КОДА ---
\lstset{
    basicstyle=\ttfamily\footnotesize,
    breaklines=true,
    texcl=true,
    frame=none
}

\DefineBibliographyStrings{russian}{
	urlseen = {дата обращения:},
}